\newpage
\subsection{Proofs using CYP syntax}
\label{sec:cyp}
Using CYP (Check Your Proof) syntax is allowed at the exams and can be a bit less to write depending on your writing style.
However, if you are very concise, you can achieve an even shorter version using a mix of CYP and non-CYP syntax.

In CYP, if doing it with pen and paper, we typically state that we use CYP for the proof,
when doing it on the computer, we need a \texttt{defs.txt} file, containing all things we assume,
as well as the function definitions, similar to Haskell syntax.
Any property we don't want to, or don't have to prove, we can denote with \texttt{axiom axiom\_name: definition of the axiom}.
Finally, we state the proof's goal, exactly as the statement to prove:
\verb|goal statement|

If we have to generalize, we add a generalized lemma and prove it, then followed by proving the specific case

For the actual proofs, it start like this:
\begin{code}{haskell}
    Lemma lemma_name: statement to prove

    Proof by induction on DataStructureHere x generalizing y -- or any other variable, or possibly without generalization
    Case Leaf -- or any other of course
        For fixed y -- only if we generalized
        Show: statement to prove for this case -- (e.g. substitute x with Leaft in this case)
        Proof
            -- The proof goes here
            statement
            (by def a_definition) .=. statement'
            (by any_axiom)        .=. statement''
        QED
    
    Case (Node x y) -- another case
        Fix x, y -- same as otherwise saying for arbitrary x, y
        Assume
            IH1: forall y: induction_hypothesis here
            IH2: forall y: another IH here -- only needed for something like this, if there are two vars.
            -- The second IH will (typically) be the same, simply a different var name
        Show: statement to prove for this case
        Proof
        -- The proof (as above)
        QED
    QED
\end{code}
